\newpage
\phantomsection
\label{prf:caratheodory-theorem}

\begin{remark*}[Return]
\hyperref[thm:caratheodory]{Return to Theorem}
\end{remark*}

\begin{theorem*}[Carathéodory's Theorem]
Let $I \subseteq \mathbb{R}$ be an open interval, let $f : I \to \mathbb{R}$, and let $c \in I$. Then $f$ is differentiable at $c$ if and only if there exists a function $\varphi : I \to \mathbb{R}$, continuous at $c$, satisfying $f(x) - f(c) = \varphi(x)(x - c)$ for all $x \in I$, in which case $\varphi(c) = f'(c)$.
\end{theorem*}

\begin{proof}
\textbf{Professional Standard Proof.}~\\
$(\Rightarrow)$ Assume $f$ is differentiable at $c$. Define $\varphi : I \to \mathbb{R}$ by
\[
  \varphi(x) :=
  \begin{cases}
    \dfrac{f(x) - f(c)}{x - c} & x \neq c, \\[6pt]
    f'(c)                       & x = c.
  \end{cases}
\]
For $x \neq c$, multiplying by $(x-c)$ gives $f(x) - f(c) = \varphi(x)(x-c)$. For $x = c$, both sides equal $0$. By definition of differentiability, $\lim_{x \to c} \varphi(x) = f'(c) = \varphi(c)$, so $\varphi$ is continuous at $c$.

$(\Leftarrow)$ Assume $\varphi : I \to \mathbb{R}$ is continuous at $c$ with $f(x) - f(c) = \varphi(x)(x-c)$ for all $x \in I$. For $x \neq c$, dividing by $(x-c)$ gives $\frac{f(x) - f(c)}{x - c} = \varphi(x)$. By continuity of $\varphi$ at $c$, $\lim_{x \to c} \frac{f(x) - f(c)}{x - c} = \varphi(c)$. Therefore $f$ is differentiable at $c$ with $f'(c) = \varphi(c)$.
\end{proof}

\begin{proof}
\textbf{Detailed Learning Proof.}~\\

\textbf{Step 1.} Establish the forward direction by constructing the required function. Assume $f$ is differentiable at $c$. Define $\varphi : I \to \mathbb{R}$ by
\[
  \varphi(x) :=
  \begin{cases}
    \dfrac{f(x) - f(c)}{x - c} & x \neq c, \\[6pt]
    f'(c)                       & x = c.
  \end{cases}
\]
This construction extends the difference quotient to all of $I$ by assigning the derivative value at $c$.

\textbf{Step 2.} Verify the factorization identity holds on all of $I$. For $x \neq c$, multiplying $\varphi(x) = \frac{f(x)-f(c)}{x-c}$ by $(x-c)$ gives $f(x) - f(c) = \varphi(x)(x-c)$ directly. For $x = c$, the left side is $f(c) - f(c) = 0$ and the right side is $\varphi(c) \cdot 0 = 0$. The identity holds on all of $I$.

\textbf{Step 3.} Show $\varphi$ is continuous at $c$. By definition of differentiability at a point,
\[
  \lim_{x \to c} \varphi(x)
  = \lim_{x \to c} \frac{f(x) - f(c)}{x - c}
  = f'(c)
  = \varphi(c).
\]
Since $\lim_{x \to c} \varphi(x) = \varphi(c)$, the function $\varphi$ is continuous at $c$.

\textbf{Step 4.} Establish the backward direction by recovering differentiability from the factorization. Assume $\varphi : I \to \mathbb{R}$ is continuous at $c$ with $f(x) - f(c) = \varphi(x)(x-c)$ for all $x \in I$. For $x \neq c$, dividing the identity by $(x - c)$ gives
\[
  \frac{f(x) - f(c)}{x - c} = \varphi(x).
\]

\textbf{Step 5.} Take the limit to establish existence of the derivative. By continuity of $\varphi$ at $c$, $\lim_{x \to c} \varphi(x) = \varphi(c)$. Therefore
\[
  \lim_{x \to c} \frac{f(x) - f(c)}{x - c} = \varphi(c).
\]

\textbf{Step 6.} Conclude differentiability and identify the derivative value. The limit of the difference quotient exists and equals $\varphi(c)$. By definition of differentiability at a point, $f$ is differentiable at $c$ with $f'(c) = \varphi(c)$.
\end{proof}

\begin{remark*}[Proof structure]
This is a bidirectional proof establishing an equivalence through explicit construction. The forward direction constructs $\varphi$ as the difference quotient extended to $c$ by the derivative value, with continuity following from the definition of the derivative. The backward direction recovers the difference quotient by algebraic manipulation and applies continuity to establish the required limit. The theorem converts the analytical condition of differentiability into the topological condition of continuity for a related function.
\end{remark*}

\begin{remark*}[Dependencies]~\\
\begin{itemize}
  \item \hyperref[def:differentiability-at-point]{Definition of Differentiability at a Point}
  \item \hyperref[def:continuity-at-point]{Definition of Continuity at a Point}
\end{itemize}
\end{remark*}